\section{Conclusion}
\label{sec:conclusion}
%%%%%%%%%%%%%%%%%%%%%%%%%%%%%%%%%%%%%%%%%%%%%%%%%%%%%%%%%%%%%%%%%%%%%%%%%%%%%%%

\subsection{What We Have Done}

This paper began with a puzzle: Why does economics remain fragmented
into separate subdisciplines---classical, behavioral, institutional, organizational---
each with its own assumptions, methods, and blind spots?

Our answer: The fragments are not competing theories but \emph{limiting cases}
of a more general structure. By identifying what each theory assumes away,
we found what unifies them.

\paragraph{The Core Innovation.}
We introduced the complementarity-context framework $(C, \Psi)$:
\begin{itemize}
  \item $C$: The complementarity matrix capturing how choices interact across agents and domains
  \item $\Psi \in \mathbb{R}^8$: The context vector decomposed into eight primitive dimensions
  \item $C^*(\Psi)$: The self-consistent reference structure---the configuration where 
        beliefs and behavior mutually reinforce
  \item $K$: The coherence index measuring alignment between actual and reference structure
  \item $Q$: The quality index measuring welfare at equilibrium
\end{itemize}

\paragraph{The Eight Context Dimensions.}
A central methodological contribution is the operationalization of context through eight independently measurable dimensions:

\begin{center}
\begin{tabular}{@{}clll@{}}
\toprule
\textbf{Symbol} & \textbf{Dimension} & \textbf{What It Captures} & \textbf{Dominant For} \\
\midrule
$\Psi_I$ & Institutional & Rules, enforcement, property rights & Democracy$\rightarrow$Growth \\
$\Psi_S$ & Social & Trust, norms, social capital & Collective action \\
$\Psi_C$ & Cognitive & Information processing, biases & Individual decisions \\
$\Psi_K$ & Informational & Transparency, signal quality & Education$\rightarrow$Earnings \\
$\Psi_E$ & Economic & Development, market thickness & Cash$\rightarrow$Consumption \\
$\Psi_T$ & Temporal & Stability, time horizons & Patience$\rightarrow$Growth \\
$\Psi_M$ & Market Scope & Geographic reach, openness & Trade effects \\
$\Psi_F$ & Factor Flexibility & Adjustment capacity, mobility & Management$\rightarrow$Productivity \\
\bottomrule
\end{tabular}
\end{center}

These dimensions emerged not from armchair theorizing but from residual analysis: we asked what systematic variation remained after standard theories had been applied.

\paragraph{Empirical Validation.}
The framework is not merely theoretical. Across six canonical economic relationships, the eight $\Psi$ dimensions explain substantial heterogeneity:

\begin{center}
\begin{tabular}{@{}lccc@{}}
\toprule
\textbf{Relationship} & \textbf{$R^2$} & \textbf{Adj. $R^2$} & \textbf{Dominant $\Psi$} \\
\midrule
Education $\rightarrow$ Earnings & 84.5\% & 76.8\% & $\Psi_K$, $\Psi_M$ \\
Management $\rightarrow$ Productivity & 78.9\% & 68.3\% & $\Psi_F$ \\
Democracy $\rightarrow$ Growth & 39.8\% & 9.7\% & $\Psi_I$, $\Psi_E$ \\
Patience $\rightarrow$ Growth & 67.3\% & 51.0\% & $\Psi_T$, $\Psi_K$ \\
Information $\rightarrow$ Behavior & 82.4\% & 73.7\% & $\Psi_K$ (--), $\Psi_F$ \\
Income $\rightarrow$ Health & 67.7\% & 51.6\% & $\Psi_T$, $\Psi_I$ \\
\midrule
\textbf{Average} & \textbf{70.1\%} & \textbf{55.2\%} & --- \\
\bottomrule
\end{tabular}
\end{center}

Three findings merit emphasis:
\begin{enumerate}
    \item \textbf{Different dimensions dominate for different effects.} Factor flexibility ($\Psi_F$) matters most for management; information quality ($\Psi_K$) for education returns; temporal stability ($\Psi_T$) for patience effects. Context is not monolithic.
    
    \item \textbf{Micro effects are better explained than macro effects.} Individual decisions show more systematic context-dependence ($R^2 \approx 75$--$85\%$) than aggregate outcomes ($R^2 \approx 40$--$70\%$). This is expected: aggregation introduces noise.
    
    \item \textbf{No ninth dimension emerged.} We tested ethnic fractionalization, natural resources, and implementation capacity. None improved the model across all outcomes. Eight dimensions appear sufficient.
\end{enumerate}

\paragraph{The Framework as Metatheory.}
Perhaps most importantly, the framework is proposed as a \emph{metatheory}---not replacing existing theories but specifying \emph{when} each applies:

\begin{itemize}
  \item \textbf{Arrow-Debreu} applies when $C_{ij} \approx 0$ and $\Psi$ is stable
  \item \textbf{Behavioral economics} applies when $\Psi_C$ varies but institutions are fixed
  \item \textbf{Institutional economics} applies when $\Psi_I$ varies but individual behavior aggregates
  \item \textbf{The full framework} applies when multiple $\Psi$ dimensions vary with significant complementarity
\end{itemize}

Existing theories tell you \emph{that} certain effects occur. Our framework tells you \emph{when} they occur.

\subsection{The Falsifiable Claim}

We commit to a specific, testable proposition:

\begin{quote}
\textbf{Core Claim:} \emph{No economic effect is context-free.}
\end{quote}

If any intervention produces identical effects across all $\Psi$ configurations---if education returns are the same in Switzerland and Nigeria, if information campaigns work equally in transparent and opaque environments, if management practices improve productivity identically with rigid and flexible labor markets---then our framework would face serious empirical challenges.

This is a strong hypothesis. Much of standard economics assumes context-independence (the ``average treatment effect'' literature). We hypothesize that context-dependence is universal and systematic.

The 70.1\% average $R^2$ provides initial support. But the claim remains falsifiable: a single context-independent effect would refute it.

\subsection{What This Changes}

\paragraph{For Economic Theory.}
The division between ``rational'' and ``behavioral'' economics is artificial. All behavior is rational in the sense of coherence-seeking. The question is not whether agents maximize utility, but what enters their utility function and how choices interact.

Classical economics is not wrong---it is the special case where $C_{ij} = 0$.
This case is useful for markets with minimal interdependence.
But it cannot explain social preferences, coordination games, or institutional dynamics.

\paragraph{For Empirical Research.}
Treatment effect heterogeneity should not be treated as noise to be averaged away---it may contain signal to be explained. The eight $\Psi$ dimensions provide a systematic framework for understanding \emph{why} effects vary across contexts.

The framework suggests a research program:
\begin{enumerate}
    \item Measure $\Psi$ dimensions for the context of interest
    \item Predict which dimensions will dominate for the effect being studied
    \item Test whether heterogeneity follows the predicted pattern
    \item Revise $\Psi$ measurement or framework structure based on results
\end{enumerate}

This is more demanding than estimating average treatment effects, but more informative.

\paragraph{For Policy Analysis.}
Standard cost-benefit analysis considers only $U^F$ (financial utility).
The FEPSDE framework shows this is radically incomplete.
Policies affect all six dimensions, and tradeoffs between dimensions may be large.

More importantly, policy effects depend on context. An intervention that works in high-$\Psi_F$ environments may fail in low-$\Psi_F$ environments. Extrapolating from one context to another requires understanding \emph{why} the original effect occurred.

\subsection{Integration with Major Research Programs}

The framework integrates with, rather than replaces, established research programs:

\begin{itemize}
    \item \textbf{Heckman}: Selection and skill formation become context-dependent (Appendix~Q)---the Heckman curve varies with $\Psi_I$ and $\Psi_{cultural}$ (Appendix~Q)
    
    \item \textbf{Bloom}: Uncertainty effects depend on (Appendix~R) $\Psi_F$ (flexibility) and $\Psi_T$ (stability)---firms respond to uncertainty differently in rigid vs. flexible environments (Appendix~R)
    
    \item \textbf{Duflo}: RCT effects are conditional (Appendix~S) on context---what works in one setting may not generalize, and $\Psi$ explains why (Appendix~S)
\end{itemize}

Each research program provides evidence for context-dependence. The framework provides the structure to unify their findings.

\subsection{Limitations Acknowledged}

We have been explicit about what the framework cannot do:
\begin{itemize}
  \item It cannot predict which equilibrium will be selected
  \item It cannot resolve normative disagreements about welfare weights
  \item It cannot model explosive dynamics when the contraction property fails
  \item Empirical validation, while promising (70.1\% $R^2$), remains preliminary
  \item The eight dimensions may have WEIRD bias from their empirical foundations
\end{itemize}

These limitations are real. But every framework has limitations.
The question is whether the framework is useful despite them---
whether it asks better questions, generates new insights, and guides research productively.

The empirical validation suggests it does.

\subsection{The Road Ahead}

\paragraph{Immediate Priorities.}
\begin{enumerate}
    \item \textbf{Replication}: Test the 8-$\Psi$ model on additional outcomes beyond the six validated here
    \item \textbf{Measurement refinement}: Develop better instruments for $\Psi_C$ (cognitive) and $\Psi_S$ (social)
    \item \textbf{Causal identification}: Move from correlational evidence to causal tests using natural experiments in $\Psi$
\end{enumerate}

\paragraph{Theoretical Extension.}
\begin{itemize}
  \item Nonlinear context dynamics with tipping points
  \item Explicit equilibrium selection mechanisms
  \item Dynamic transition analysis between $\Psi$ configurations
  \item Strategic context manipulation by large actors
\end{itemize}

\paragraph{Policy Application.}
\begin{itemize}
  \item Trade policy design that accounts for identity effects and $\Psi$-heterogeneity
  \item Climate governance as a global coherence problem across nations with different $\Psi$
  \item Development interventions calibrated to local $\Psi$ configurations
  \item Institutional reform that avoids coherence traps
\end{itemize}

\subsection{Final Reflection}

Economics began as moral philosophy---the study of how societies organize
to meet human needs. Over time, it narrowed to the study of market efficiency.
Behavioral economics recovered some of the lost terrain.
Institutional economics recovered more.

This framework attempts to reunify the fragments.
It says: Economics is the study of coherence---
how beliefs, behaviors, and contexts align (or fail to align)
to create stable patterns that determine human welfare.

The eight $\Psi$ dimensions operationalize ``context'' so it can be measured, tested, and used.
The 70.1\% explanatory power shows this operationalization has empirical bite.
The falsifiable core claim---no effect is context-free---distinguishes the framework from unfalsifiable alternatives.

Efficiency is one aspect of coherence, the limiting case where choices don't interact.
But most of economic life involves choices that do interact---
creating complementarities, externalities, coordination problems, and social dilemmas.

Understanding these interactions requires a richer framework than
``maximize utility subject to constraints.''
It requires understanding what enters utility (FEPSDE),
how utilities interact ($C$), how context shapes interactions ($\Psi$),
and how all three co-evolve.

The Complementarity-Context framework provides this richer structure.
It is not complete. It is not final. But it is, we hope, a step toward
a more unified and more useful economics---
one that can address the great challenges of our time:
climate change, inequality \citep{piketty2014}, institutional decay, technological disruption,
and the search for coherent ways of living together on a finite planet.

\vspace{1em}
\begin{center}
\emph{Rationality is not maximization. It is coherence.} \\
\emph{Welfare is not wealth. It is flourishing across dimensions.} \\
\emph{Context is not noise. It is the key to generalization.} \\[0.5em]
\emph{And with eight measurable dimensions, we can finally put this vision to the test.}
\end{center}
